\section{Discussion and Limitations}
\label{sec:discussion}
\subsection{Scope of Recovery Guarantees}
The recovery gate exhaustively evaluates a finite deterministic scenario set.
Within this model, a feasible output satisfies the stated deadline and
recovery limits. This statement should not be extended to arbitrary failure
times, simultaneous faults, all replica-transfer failures, or stochastic
real-time executions. Node and copy failures are injected at different
predefined times, and the replay is not a general event-driven model of
every possible execution history. Broader failure coverage would require an
expanded scenario set and a corresponding re-evaluation.

Domain-separated replicas provide a construction-level protection rule, but
the default gate tests single-copy, node, and effective-transfer faults.
Consequently, an unchanged single-point feasible rate after relaxing domain
separation would not establish that separation is unnecessary under
correlated outages. Conversely, correlated-outage results must specify which
domains fail and whether dynamic reexecution remains available. We do not
claim a domain-failure advantage without that additional experiment.

\subsection{Learning Signal and Generalization}
The scorer imitates feasible, low-cost teacher deployments rather than
globally optimal assignments. A candidate not chosen by the teacher can
still participate in a good deployment. Limited assignment-label coverage,
limited supervised validation instances, and synthetic feature distributions
can therefore constrain generalization.

The resource-graph score table is deployment-independent and fixed during a
search. This makes inference inexpensive relative to repeatedly evaluating a
network, but it cannot explicitly react to changes in partial resource
occupancy or to the evolving recovery bottleneck. The feasibility-first
objective, repair logic, and scenario gate continue to carry the main
constraint-handling responsibilities. End-to-end learning utility must be
established empirically rather than inferred from classification AUC.

\subsection{Search and Cost Interpretation}
The recovery surrogate is not an admissible bound. Screening on that
surrogate can prevent some truly feasible deployments from being checked.
The distinction between the quick deployment cost and the full
recovery-augmented objective also prevents a claim of monotonic full-cost
improvement. The solver should be viewed as a bounded heuristic returning
a scenario-verified feasible deployment when it finds one.

Communication charges cover all provisioned replica pairs while readiness
selects an effective producer. Resource limits constrain cumulative deployed
footprint rather than instantaneous CPU sharing, and the recovery charge is
a uniformly averaged delay proxy. These are explicit modeling choices;
different charging or reservation models can change the preferred
replication policy. Failure-rate features guide placement but do not turn
the uniform scenario average into a probability-weighted risk measure.

\subsection{External Validity and Baseline Adaptation}
The custom workflow shapes and synthetic machine snapshots support
controlled comparisons but do not constitute validation on real executions.
Task CPU and memory demands are homogeneous in the present synthetic path,
and fault domains are generated rather than measured. Real workflow traces,
measured resource demands, and observed correlated-failure groups are
important follow-up evaluations.

The external baselines also face an adaptation limitation. Mapping a
multi-objective reliability algorithm or an uncertainty-aware hybrid method
onto deterministic, fixed-order, fixed-resource scheduling can remove some
of its strengths. A shared verifier supports consistent scoring of the
adapted deployments; it does not eliminate that limitation. The strongest
method-specific comparison for the learned increment is therefore the
nonlearned ALNS ablation under the same search and recovery framework.
